
\subsection{Algorithmes de jointure}

\begin{frame}[fragile]{L'opérateur qui nous manque}

La jointure: opération très courante, potentiellement coûteuse.

\vfill

L'opérateur de jointure complète notre petit catalogue pour (presque) toutes 
les requêtes SQL.

\vfill

\begin{minted}[frame=leftline,
               baselinestretch=.9,
               fontsize=\footnotesize,
               framesep=2mm]{sql}
    select a1, a2, .., an
    from T1, T2, ..., Tm
    where T1.x = T2.y and ... 
    order by ...
\end{minted}

\vfill

Plusieurs algorithmes possibles.
\end{frame}

\begin{frame}{Principaux algorithmes}

\alert{Jointure avec index}

\begin{itemize}

  \item Avec un index\,: \alert{jointure par boucles imbriquées indexées}.
   
  \item  Avec deux index\,: on fait comme si on avait
            un seul index
\end{itemize}

\vfill

\alert{Jointure sans index}

\begin{itemize}

  \item Le plus simple\,: \alert{jointure
           par boucles imbriquées} (non indexée)

  \item Le plus courant\,: \alert{jointure par tri-fusion} 

   \item Parfois le meilleur\,: \alert{jointure par hachage}
\end{itemize}
\end{frame}

%%%%%%%%%%%%%%%%%%%%%%%%%%%%%%%%%%%%%%%%%%%%%%%%%%%%%%%
\begin{frame}[fragile]{Jointure avec index}

C'est la jointure la plus \alert{courante} car on effectue
\alert{naturellement} la jointure sur les clés primaires/étrangères.

\begin{itemize}

 \item Les films et leur metteur en scène
 
\begin{minted}[frame=leftline,
               baselinestretch=.9,
               fontsize=\footnotesize,
               framesep=2mm]{sql}
      select * from Film as f, Artiste as a 
      where f.id_realisateur = a.id
\end{minted}
      
 \item Les artistes et leurs rôles      
 
 \begin{minted}[frame=leftline,
               baselinestretch=.9,
               fontsize=\footnotesize,
               framesep=2mm]{sql}
      select * from Artiste as a, Role as r
      where a.id = r.id_acteur
\end{minted}

      
 \item Les employés et leur département

\begin{minted}[frame=leftline,
               baselinestretch=.9,
               fontsize=\footnotesize,
               framesep=2mm]{sql}
      select * from emp e, dept d
      where e.dnum = d.num
\end{minted}

\end{itemize}

Garantit \alert{au moins} un index sur la condition de jointure.
\end{frame}

\begin{frame}{Jointure avec index: l'algorithme}

Algorithme dit de \alert{boucles imbriquées indexées}.

\begin{itemize}

  \item on parcourt séquentiellement la table contenant la clé étrangère;
 
   \item pour chaque tuple, on utilise la valeur de la clé étrangère pour
   accéder à l'index sur la clé primaire
   de la second table: on récupère l'adresse \texttt{adr} d'un tuple;
 \item il reste à effectuer un accès direct, avec l'adresse \texttt{adr}, pour
   obtenir le second tuple et constituer la paire. 

\end{itemize}

\vfill

Avantages\,:
  \begin{itemize}
      \item  Efficace (un parcours, plus des recherches par adresse)

      \item Favorise le temps de réponse \textbf{et} le temps
                 d'exécution
\end{itemize}

\end{frame}

\begin{frame}{L'opérateur}

Réutilise des méthodes existantes.

\vfill

\figSlideWithSize{planEx-indexedjoin}{5cm}

\end{frame}

\begin{frame}[fragile]{Le code? Simplissime}
\begin{minted}[frame=leftline,
               baselinestretch=.9,
               fontsize=\footnotesize,
               framesep=2mm]{bash}
    function nextIndexJoin
    {
      # $tScan est le parcours de la table      
      $tuple = $tScan.next();

      # On instancie un parcours d'index
      iScan = new IndexScan ();
      # Le parcours d'index avec la cle etrangere
      $iScan.open ($tuple.foreignKey);
      # On prend l'adresse 
      addr = $iScan.next();
 
      # Et on renvoie la paire avec le tuple et l'adresse      
      return [$tuple, $addr];
     }
\end{minted}

\end{frame}

%%%%%%%%%%%%%%%%%%%%%%%%%%%%%%%%%%%%%%%%%%%%%%%%%%%%%%%
\begin{frame}{Et avec deux index\,?}

On pourrait penser à la solution suivante\,:

\begin{itemize}
  \item Fusionner les deux index\,: on obtient des paires
                d'adresse.

  \item Pour chaque paire, aller chercher la ligne $A$,
         la ligne $B$.
\end{itemize}

Problématique car beaucoup d'accès aléatoires
(cf. \alert{Quand utiliser un index}). 

\vfill

En pratique\,:

  \begin{itemize}
      \item On se ramène à la jointure avec un index

      \item On prend la petite table comme table extérieure.
\end{itemize}
\end{frame}


%%%%%%%%%%%%%%%%%%%%%%%%%%%%%%%%%%%%%%%%%%%%%%%%%%%%%%%

\begin{frame}[fragile]{Jointures par boucles imbriquées}

Pas d'index\,? On fait bête et méchant: énumération de toutes
les solutions.

\begin{minted}[frame=leftline,
               baselinestretch=.9,
               fontsize=\footnotesize,
               framesep=2mm]{bash}
    function JoinList
    {
      # $L1, $L2, listes à joindre
      # $condition est la condition de jointure
      resultat = [];
      for tuple1 in $L1
      do
       for tuple2 in $L2
       do
         if (condition ($tuple1, $tuple2) = true) the
           $resultats[] = ($tuple1, $tuple2);
         fi
       done
     done
   }
\end{minted}

C'est la \alert{jointure par boucles imbriquées}
\end{frame}


\begin{frame}{Comment parcourir les tables}

On prend toutes les paires de blocs, on les joint.
\vfill

\figSlideWithSize{nestedloop}{12cm}

\vfill

Peu de mémoire, mais coût très élevé: $B_R + B_R \times B_S$.

\end{frame}
%%%%%%%%%%%%%%%%%%%%%%%%%%%%%%%%%%%%%%%%%%%%%%%%%%%%%%%
\begin{frame}{Essentiel\,: la mémoire}

On alloue le maximum à la table intérieure de la boucle
imbriquée.

\medskip


\figSlideWithSize{nestedloop-improved}{12cm}

\medskip

Si la table intérieure tient en mémoire\,: une seule
lecture des deux tables suffit: $B_R + B_S$.

\end{frame}
%%%%%%%%%%%%%%%%%%%%%%%%%%%%%%%%%%%%%%%%%%%%%%%%%%%%%%%
\begin{frame}{Jointure par tri fusion}

Plus efficace que les boucles imbriquées pour de grosses
tables.

\begin{itemize}

  \item On trie les deux tables sur les colonnes de jointures

  \item On effectue la fusion
\end{itemize}

\vfill

C'est le tri qui coûte cher.

\vfill
\begin{block}{Important}
\alert{Opérateur bloquant:}\,: on ne peut {\em rien} obtenir
tant que le tri n'est pas fini.
\end{block}
\end{frame}
%%%%%%%%%%%%%%%%%%%%%%%%%%%%%%%%%%%%%%%%%%%%%%%%%%%%%%%
\begin{frame}[fragile]{La procédure de fusion}

On tri $R$ et $S$ et on parcourt ensuite les tables
triées en parallèle.

\vfill

La fusion:on commence avec les premiers enregistrements $r_1$ 
et $s_1$  de chaque 
table. 

\begin{itemize}
   \item Si $r_1.a = s_1.b$, on joint les deux enregistrements,  on passe au
      enregistrements suivants, jusqu'à ce que $r_i.a \not= s_i.b$.
    \item Si $r_1.a < s_1.b$, on avance sur la liste de $R$. 
    \item Si $r_1.a > s_1.b$, on avance sur la liste de $S$. 
\end{itemize}


\end{frame}


\begin{frame}{Tri-fusion\,: illustration}

\figSlideWithSize{sortmerge}{7.5cm}

\end{frame}


\begin{frame}{Tri-fusion\,: l'opérateur}

On a déjà le tri: un opérateur de fusion (\texttt{Merge}) vient le compléter.


\figSlideWithSize{planEx-trifusion}{7.5cm}

\end{frame}

%%%%%%%%%%%%%%%%%%%%%%%%%%%%%%%%%%%%%%%%%%%%%%%%%%%%%%%
\begin{frame}{Jointure par hachage}

En théorie le plus efficace.

\begin{itemize}

  \item Très rapide quand
         une des deux tables est petite (1, 2, 3 fois
           la taille de la mémoire).

   \item Pas très robuste (efficacité dépend de plusieurs
      facteurs).
\end{itemize}

\vfill

Principe:
\begin{itemize}

  \item On hache la plus petite des deux tables en $n$ fragments.

  \item On hache la seconde table, avec la même fonction $h()$, 
             en $n$ autres fragments.

  \item On réunit les fragments par paire, et on fait la jointure.
\end{itemize}

\alert{Essentiel}\,: pour chaque paire, au moins un fragment
doit tenir en mémoire.

\end{frame}
%%%%%%%%%%%%%%%%%%%%%%%%%%%%%%%%%%%%%%%%%%%%%%%%%%%%%%%
\begin{frame}{Illustration\,: phase de hachage}

On suppose que $R$ est la plus petite.  

\vfill

On détermine $k$ pour que
$\frac{|R|}{k} < M$, la taille de la mémoire.

\vfill

\figSlideWithSize{hashjoin}{9cm}

\vfill

On hache $S$ en $k$ fragments avec la même fonction.

\end{frame}

\begin{frame}{Jointure: la propriété clé}

Un enregistrement  $r$ de $R$ est  
placé dans le fragment  $F_R^{h(r.a)}$;

\vfill

un enregistrement  $s$ de $S$ est  
placé dans le fragment  $F_S^{h(s.b)}$.

\vfill

\begin{block}{Propriété}
Si deux tuples  $r$ et  $s$
doivent être joints, alors on a  $h(r.a) = h(s.b)=u$ 
et on les trouvera, respectivement, dans  $F_R^u$
et $F_S^u$.
\end{block}

\vfill

\alert{En d'autres termes, il suffit d'effectuer
la jointure sur les paires de fragments correspondant à la même valeur
de la fonction de hachage.}
\end{frame}


%%%%%%%%%%%%%%%%%%%%%%%%%%%%%%%%%%%%%%%%%%%%%%%%%%%%%%%
\begin{frame}{Illustration\,: phase de jointure}

On chaque un fragment $F_R^i$ de $R$ en mémoire (il tient).

\vfill

On parcourt le fragment $F_S^i$ de $S$ et on joint.

\vfill

\figSlideWithSize{hashjoin2}{9cm}

\vfill

\alert{Déjà vu?} Oui: on s'est ramené à la jointure par boucles imbriquées
quand une table tient en mémoire.
\end{frame}

\begin{frame}{L'opérateur de jointure par hachage}
Comment implanter cet algorithme de jointure sous forme d'itérateur? 

\begin{itemize}
  \item Toute la phase de hachage s'effectue dans le \texttt{open()}
  \item Le \texttt{next()} avance d'une étape dans la jointure de deux fragments.
\end{itemize}

\alert{Cet opérateur est donc bloquant}

\end{frame}


%%%%%%%%%%%%%%%%%%%%%%%%%%%%%%%%%%%%%%%%%%%%%%%%%%%%%%%
\begin{frame}{À retenir}

Pour des requêtes effectuées régulièrement: il faut utiliser
la jointure par boucle imbriquées indexée.

\begin{itemize}
  \item si nécessaire, créer un index;
  \item si jointure sur clés, l'index doit \alert{déjà} exister.
\end{itemize}

\vfill

Si jointure sur attributs non indexés
\begin{itemize}
  \item l'optimiseur est censé prendre la bonne décision;
  \item ça risque d'être long si aucune des deux tables ne tient en mémoire.
\end{itemize}

\vfill

Requête très lente? \alert{Vérifier l'algorithme choisi pour la jointure.}
\end{frame}